% Agent 26: Pages 443--446 (PDF pages 50--53)
% Sections 5.11, 5.12, 5.13, 5.14, 6, 6.1

%%%%%%%%%%%%%%%%%%%%%%%%%%%%%%%%%%%%%%%%%%%%%%%%%%%%%%%%%%%%%%%%%%%
%% Figure 5.10.1 (appears on book page 442)
%%%%%%%%%%%%%%%%%%%%%%%%%%%%%%%%%%%%%%%%%%%%%%%%%%%%%%%%%%%%%%%%%%%

\begin{figure}[htbp]
\centering
% \includegraphics[width=\columnwidth]{fig_5_10_1}
\fbox{\parbox{0.9\columnwidth}{\centering [Figure 5.10.1 placeholder]\\
The total power per unit frequency $L_\nu$ emitted by two model disks around
black holes, as calculated by Shakura and Sunyaev~(1972) without taking account of
capture of radiation by the hole. In both models the hole was
assumed to be nonrotating, so the inner edge of the disk was at $r_* = r/h = 6$. Model~(a)
corresponds to
$\alpha \sim 10^{-3}$, \quad $M = M_\odot$, \quad $M = 10^{38}\;\text{erg/sec}$;\\
$L \simeq L_{\text{crit}} = 10^{38}\;\text{erg/sec}$;\\
$\dot{M}_0 = 10^{-9}\,M_\odot/\text{yr}$, \quad $L = 10^{36}\;\text{erg/sec}$.\\
Model~(b) corresponds to
$\alpha \sim 10^{-2}$ to $1$ (spectrum insensitive to $\alpha$),\\
$M = M_\odot$, \quad $\dot{M}_0 = 10^{-9}\,M_\odot/\text{yr}$, \quad $L = 10^{36}\;\text{erg/sec}$.}}
\caption{The total power per unit frequency $L_\nu$ emitted by two model disks around
black holes. See text for parameters.}
\label{fig:5.10.1}
\end{figure}

%%%%%%%%%%%%%%%%%%%%%%%%%%%%%%%%%%%%%%%%%%%%%%%%%%%%%%%%%%%%%%%%%%%
%% §5.11 Heating of the Outer Region by X-Rays from the Inner Region
%%%%%%%%%%%%%%%%%%%%%%%%%%%%%%%%%%%%%%%%%%%%%%%%%%%%%%%%%%%%%%%%%%%

\subsection{Heating of the Outer Region by X-Rays from the Inner Region}
\label{sec:5.11}

Thus far we have ignored any radiative energy transfer from one portion of the
disk to another. As long as the disk is optically thick, no significant transfer can
occur through its interior. However, radiation can skim along the surface of the
disk from the hot, inner region toward the cooler outer region. Because the disk
is thick with increasing radius
%
\[
h = \text{constant} \quad \text{in inner region},
\]
\[
h \propto r^{21/20} \quad \text{in middle region},
\]
\[
h \propto r^{9/8} \quad \text{in outer region},
\]
%
such radiation will get captured in the middle and outer regions and will heat
their surface layers. As a result, the outer regions will thicken even more than is
predicted by the above model. Rough estimates by Shakura and Sunyaev~(1972)
\cite{ShakuraSunyaev1972}
suggest that for typical binary accretion models ${\sim}\,1$ per cent of the total X-ray
luminosity can be captured and reradiated. About 10 per cent of the reradiated
flux comes off in spectral lines in the optical region, and about 90 per cent comes
off in ultraviolet free-bound and free-free emission.

%%%%%%%%%%%%%%%%%%%%%%%%%%%%%%%%%%%%%%%%%%%%%%%%%%%%%%%%%%%%%%%%%%%
%% §5.12 Fluctuations on the Steady-State Model
%%%%%%%%%%%%%%%%%%%%%%%%%%%%%%%%%%%%%%%%%%%%%%%%%%%%%%%%%%%%%%%%%%%

\subsection{Fluctuations on the Steady-State Model}
\label{sec:5.12}

Superimposed on the steady-state disk structure will be local fluctuations due to
turbulence, magnetic flux reconnection, and various plasma instabilities. Such
fluctuations will probably be large if $\alpha \sim 1$, and unimportant if $\alpha \ll 1$.

Local fluctuations will presumably create local, extra-luminous hot spots.
Sunyaev~(1972) \cite{Sunyaev1972} estimates that such hot spots can last for many orbital periods.
Doppler shifts and focusing by the black hole will cause the frequency and
intensity of the radiation from a hot spot as seen at Earth to fluctuate with the
orbital period. The minimum possible fluctuation period is the orbital period at
the last stable circular orbit:

\begin{equation}
P_{\min} = 12\pi \sqrt{6}\, GM/c^3 = (0.5 \;\text{msec})(M/M_\odot)
\quad \text{for nonrotating hole,}
\tag{5.12.1a}
\end{equation}
\begin{equation}
P_{\min} = 4\pi M = (0.06 \;\text{msec})(M/M_\odot)
\quad \text{for maximally rotating hole.}
\tag{5.12.1b}
\end{equation}

In a binary system one can measure with some confidence the mass of the
compact, accreting object by examining the Doppler shifts of spectral lines
(lines from the normal star and from the outer regions of the disk, not from
the inner disk.) If one could also find quasiperiodic fluctuations in the X-ray
output and place a lower limit on their periods, one might be able to distinguish
whether the black hole is nonrotating ($P_{\min} \sim (0.5\;\text{msec})(M/M_\odot)$) or rotating
$[(0.06\;\text{msec})(M/M_\odot) \leq P_{\min} \leq (0.5\;\text{msec})(M/M_\odot)]$. This ``test for rotation'' was
devised by the Moscow group of Zel'dovich, Novikov, Polnarev, and Sunyaev
(see Sunyaev 1972) \cite{Sunyaev1972}.

An alternative test for rotation devised by Thorne~(1973b) \cite{Thorne1973b} involves the shape
of Doppler-broadened spectral features. The ``dragging of inertial frames'' by the
rotation of the hole will produce a strong asymmetry in the Doppler broadening
of any spectral features emitted from in or near the ergosphere. Three factors
make such asymmetries a less attractive test for rotation than the Moscow test:
(i) the spectrum from in and near the ergosphere should be rather featureless
because of the high temperatures; (ii) the Doppler broadening is so great that
it would smear into the background all features except abnormally strong ones;
(iii) Comptonization can also produce asymmetries in spectral features.

Local fluctuations (``hot spots'') may well produce nonthermal radiation. A
case of particular interest is the synchrotron radiation emitted by charged particles
that are accelerated in regions of reconnecting magnetic flux (see \S2.9). For
supermassive holes at the centers of galaxies, Lynden-Bell~(1969) \cite{LyndenBell1969} has estimated
the spectrum of such synchrotron radiation. He argues that the mean free paths
of electrons against Coulomb scattering will be too small to permit acceleration
in the reconnection regions; but that protons, having larger inertia, can get
accelerated significantly. The result is a large amount of proton synchrotron
radiation in the radio band. Efforts are being made to compare the predictions
of this model with radio-frequency observations of the galactic center (Ekers
and Lynden-Bell 1970) \cite{EkersLyndenBell1970}.

%%%%%%%%%%%%%%%%%%%%%%%%%%%%%%%%%%%%%%%%%%%%%%%%%%%%%%%%%%%%%%%%%%%
%% §5.13 Supercritical Accretion
%%%%%%%%%%%%%%%%%%%%%%%%%%%%%%%%%%%%%%%%%%%%%%%%%%%%%%%%%%%%%%%%%%%

\subsection{Supercritical Accretion}
\label{sec:5.13}

If the accretion rate onto a black hole were larger than the ``Eddington limit'',

\begin{equation}
\dot{M}_{\text{crit}} \sim (1 \times 10^{-8}\, M_\odot/\text{yr})(M/M_\odot),
\tag{5.13.1}
\end{equation}

the luminosity produced by our model would exceed the ``Eddington limit'',

\begin{equation}
L_{\text{crit}} = (1 \times 10^{38}\;\text{erg/sec})(M/M_\odot).
\tag{5.13.2}
\end{equation}

[See eq.~(4.5.4) and \S5.1.] Shakura and Sunyaev~(1972) \cite{ShakuraSunyaev1972} suggest the following
general picture for disk accretion in this ``supercritical'' case: In the outer regions
of the disk the gas is shielded from radiation pressure by high opacity, so it
accretes in the usual manner. However, as the gas nears the black hole, radiation
pressure builds up, and ultimately becomes strong enough to eject the gas out of
the disk, in the $z$-direction. Only a small fraction of the accreting matter ever
reaches the hole, and the total luminosity is self-regulated at the Eddington
critical value (\S5.1.2).

Rough computations by Shakura and Sunyaev suggest that, if the accretion
is strongly supercritical

\begin{equation}
\dot{M}_0 \gg \dot{M}_{\text{crit}} \times (10^3 a_* M_\odot / M)^{2/3},
\tag{5.13.3}
\end{equation}

the outflowing gas becomes opaque and reprocesses the radiation emitted near
the hole from high frequencies to lower frequencies. Most of the energy comes
off in the ultraviolet and optical regions; and the outflowing matter achieves
velocities

\begin{equation}
v \sim (10^5 \;\text{cm/sec})\, \alpha (\dot{M}_0 / \dot{M}_{\text{crit}}).
\tag{5.13.4}
\end{equation}

The opaque region in this case is ${\sim}\,10^{10}$ to $10^{12}$~cm in radius, and may thus
cover the normal star as well as the disk and the hole.

%%%%%%%%%%%%%%%%%%%%%%%%%%%%%%%%%%%%%%%%%%%%%%%%%%%%%%%%%%%%%%%%%%%
%% §5.14 Comparison with Observations
%%%%%%%%%%%%%%%%%%%%%%%%%%%%%%%%%%%%%%%%%%%%%%%%%%%%%%%%%%%%%%%%%%%

\subsection{Comparison with Observations}
\label{sec:5.14}

Currently the time scale for 100 per cent improvements in the observations of
X-ray sources is less than a year. Therefore, it would be foolhardy to present in
these notes a detailed comparison of the above models with X-ray observations.
We shall merely remark that as of August 1972 the comparisons are very
promising. In particular, the X-ray sources Cyg X-1 and 2U\,0900-40 are excellent
binary black-hole candidates.

The comparison of these models with observations of galactic nuclei is much
more difficult than comparison with X-ray source observations. Obscuration by
dust, and the high density of radiation sources in galactic nuclei create great
difficulties. However, the observations are not in conflict with the hypothesis
that a supermassive hole of $M \lesssim 10^9\,M_\odot$ resides at the center of our Galaxy. (See
Lynden-Bell and Rees 1971.) \cite{LyndenBellRees1971}

%%%%%%%%%%%%%%%%%%%%%%%%%%%%%%%%%%%%%%%%%%%%%%%%%%%%%%%%%%%%%%%%%%%
%% §6 White Holes and Black Holes of Cosmological Origin
%%%%%%%%%%%%%%%%%%%%%%%%%%%%%%%%%%%%%%%%%%%%%%%%%%%%%%%%%%%%%%%%%%%

\section{White Holes and Black Holes of Cosmological Origin}
\label{sec:6}

%%%%%%%%%%%%%%%%%%%%%%%%%%%%%%%%%%%%%%%%%%%%%%%%%%%%%%%%%%%%%%%%%%%
%% §6.1 White Holes, Grey Holes, and Black Holes
%%%%%%%%%%%%%%%%%%%%%%%%%%%%%%%%%%%%%%%%%%%%%%%%%%%%%%%%%%%%%%%%%%%

\subsection{White Holes, Grey Holes, and Black Holes}
\label{sec:6.1}

All black holes studied in previous sections were regarded as remnants of the
relativistic collapse of massive stars. However, the formation of stars from rarefied
gas is possible only in relatively late stages of the evolution of the Universe
($t \gtrsim 10^7$~years). In earlier stages, according to current cosmological models,
radiation pressure impeded the growth of condensations; so individual bodies
could not form from \emph{small} perturbations of the primeval gas. In the absence of
star formation, there correspondingly should have been no black-hole formation.
However, this conclusion might be wrong. It is quite possible that near the
beginning of the cosmological expansion the matter distribution and metric were
\emph{highly} inhomogeneous and anisotropic. In this case the condensation of individual
bodies would have been possible. Moreover, because in the early stages of the hot
universe the mass-energy in radiation and in particle-antiparticle pairs greatly
exceeded that in baryons, bodies which condensed then must have been made
primarily of photons and pairs. However, because such an ultrarelativistic gas has
an adiabatic index $\Gamma \leq 4/3$, any body made from it is unstable against gravita\-
tional collapse. Thus, one is led to imagine primordial condensations that, like
Wheeler's~(1955) \cite{Wheeler1955} geons, were made primarily from photons and pairs, but unlike
geons, quickly collapsed to form black holes.

Another possible type of strong inhomogeneity in the early universe is a
``white hole'' (Novikov 1964; Ne'eman 1965) \cite{Novikov1964, Neeman1965}. Suppose that some isolated
portions of the universe failed to begin their expansion at the same moment as
the rest of the universe. As measured by external observers, the delay in the
initial expansion of a given region (or ``core'') might be arbitrarily great; and the
delays might vary from core to core. When a ``lagging core'' eventually begins to
expand, and emerges through its gravitational radius, external observers should
see an explosion with a release of tremendous energy. Such an object is, in some
sense, a concrete realization of Ambartsumyan's (1961, 1964) concept of a super\-
dense ``D-body''. (Thus, one can regard ``D-body'', ``white hole'', and ``lagging
core'' as different names for the same concept.)

Finally, if a lagging core, when it begins to expand, has insufficient energy to
emerge through its gravitational radius into the external universe, then one can
regard it as a ``grey hole''. Thus, in the idealized spherical case, the matter of a
grey hole is ``born'' in the Kruskal diagram, staying always to the left of the
center line so it is always separated from the external universe by a past horizon
or a future horizon; and it finally dives to its death in the terminal $r = 0$
singularity.

Might all three types of cosmological holes actually exist in Nature? What
would be their properties? Theorists have very little observational way to rule
out their existence; and they have said very little about them.

In these notes we shall dwell on only two points: (i) the growth of cosmological
holes by accretion, and (ii) a limit on how many baryons can be inside cosmological
holes.
